% Generated by tools/build_prior_work.py from the claims ledger; do not edit.
\bibitem{pwadams2003sobolev}
Adams, R. A., and Fournier, J. J. F. (2003). \emph{Sobolev Spaces}. Academic Press. ISBN 978-0-12-044143-3.

\bibitem{pwaghanim2020planck}
Aghanim, N., and others (2020). ``Planck 2018 results. VI. Cosmological parameters.'' \emph{Astronomy \& Astrophysics} \textbf{641}, A6. \url{https://doi.org/10.1051/0004-6361/201833910}

\bibitem{pwaizenman1981third}
Aizenman, M., and Lieb, E. H. (1981). ``The third law of thermodynamics and the degeneracy of the ground state for lattice systems.'' \emph{Journal of Statistical Physics} \textbf{24}, 279--297. \url{https://doi.org/10.1007/BF01007649}

\bibitem{pwalicki2007quantum}
Alicki, R., and Lendi, K. (2007). \emph{Quantum Dynamical Semigroups and Applications}. Springer. \url{https://doi.org/10.1007/3-540-70861-8}

\bibitem{pwarnold1989mathematical}
Arnold, V. I. (1989). \emph{Mathematical Methods of Classical Mechanics}. Springer, New York (Graduate Texts in Mathematics). \url{https://doi.org/10.1007/978-1-4757-2063-1}

\bibitem{pwarnold2006finite}
Arnold, D. N., Falk, R. S., and Winther, R. (2006). ``Finite element exterior calculus, homological techniques, and applications.'' \emph{Acta Numerica} \textbf{15}, 1--155. \url{https://doi.org/10.1017/S0962492906210018}

\bibitem{pwbaader1998term}
Baader, F., and Nipkow, T. (1998). \emph{Term Rewriting and All That}. Cambridge University Press. \url{https://doi.org/10.1017/CBO9781139172752}

\bibitem{pwbarbour1982mach}
Barbour, J. B., and Bertotti, B. (1982). ``Mach's principle and the structure of dynamical theories.'' \emph{Proceedings of the Royal Society of London. A. Mathematical and Physical Sciences} \textbf{382}, 295--306. \url{https://doi.org/10.1098/rspa.1982.0102}

\bibitem{pwbarnum1996noncommuting}
Barnum, H., Caves, C. M., Fuchs, C. A., Jozsa, R., and Schumacher, B. (1996). ``Noncommuting mixed states cannot be broadcast.'' \emph{Physical Review Letters} \textbf{76}, 2818--2821. \url{https://doi.org/10.1103/PhysRevLett.76.2818}

\bibitem{pwbipm2019international}
Bureau International des Poids et Mesures (2019). ``The International System of Units (SI).'' \url{https://www.bipm.org/en/publications/si-brochure}

\bibitem{pwboffi2010finite}
Boffi, D. (2010). ``Finite element approximation of eigenvalue problems.'' \emph{Acta Numerica} \textbf{19}, 1--120. \url{https://doi.org/10.1017/S0962492910000012}

\bibitem{pwbossavit1988whitney}
Bossavit, A. (1988). ``Whitney forms: a class of finite elements for three-dimensional computations in electromagnetism.'' \emph{IEE Proceedings A: Physical Science, Measurement and Instrumentation, Management and Education, Reviews} \textbf{135}, 493--500. \url{https://doi.org/10.1049/ip-a-1.1988.0077}

\bibitem{pwbratteli1972inductive}
Bratteli, O. (1972). ``Inductive limits of finite dimensional C*-algebras.'' \emph{Transactions of the American Mathematical Society} \textbf{171}, 195--234. \url{https://doi.org/10.1090/S0002-9947-1972-0312282-2}

\bibitem{pwbratteli1987operator}
Bratteli, O., and Robinson, D. W. (1987). \emph{Operator Algebras and Quantum Statistical Mechanics 1: C*- and W*-Algebras, Symmetry Groups, Decomposition of States}. Springer. \url{https://doi.org/10.1007/978-3-662-02520-8}

\bibitem{pwbratteli1997operator}
Bratteli, O., and Robinson, D. W. (1997). \emph{Operator Algebras and Quantum Statistical Mechanics 2}. Springer. \url{https://doi.org/10.1007/978-3-662-03444-6}

\bibitem{pwbrenner2008mathematical}
Brenner, S. C., and Scott, L. R. (2008). \emph{The Mathematical Theory of Finite Element Methods}. Springer, New York (Texts in Applied Mathematics). \url{https://doi.org/10.1007/978-0-387-75934-0}

\bibitem{pwbrink1976local}
Brink, L., Deser, S., Zumino, B., Di Vecchia, P., and Howe, P. (1976). ``Local supersymmetry for spinning particles.'' \emph{Physics Letters B} \textbf{64}, 435--438. \url{https://doi.org/10.1016/0370-2693(76)90115-5}

\bibitem{pwbrocker1985representations}
Br{\"{o}}cker, T., and tom Dieck, T. (1985). \emph{Representations of Compact Lie Groups}. Springer. \url{https://doi.org/10.1007/978-3-662-12918-0}

\bibitem{pwbrouwer2012spectra}
Brouwer, A. E., and Haemers, W. H. (2012). \emph{Spectra of Graphs}. Springer, New York (Universitext). \url{https://doi.org/10.1007/978-1-4614-1939-6}

\bibitem{pwbroy2001specification}
Broy, M., and St{\o}len, K. (2001). \emph{Specification and Development of Interactive Systems}. Springer. \url{https://doi.org/10.1007/978-1-4613-0091-5}

\bibitem{pwburnside1905condition}
Burnside, W. (1905). ``On the condition of reducibility of any group of linear substitutions.'' \emph{Proceedings of the London Mathematical Society} \textbf{s2-3}, 430--434. \url{https://doi.org/10.1112/plms/s2-3.1.430}

\bibitem{pwbusch1991informationally}
Busch, P. (1991). ``Informationally complete sets of physical quantities.'' \emph{International Journal of Theoretical Physics} \textbf{30}, 1217--1227. \url{https://doi.org/10.1007/BF00671008}

\bibitem{pwbusch2003quantum}
Busch, P. (2003). ``Quantum states and generalized observables: a simple proof of Gleason's theorem.'' \emph{Physical Review Letters} \textbf{91}, 120403. \url{https://doi.org/10.1103/PhysRevLett.91.120403}

\bibitem{pwbyckling1973particle}
Byckling, E., and Kajantie, K. (1973). \emph{Particle Kinematics}. Wiley. ISBN 0-471-12885-6.

\bibitem{pwchiribella2011informational}
Chiribella, G., D'Ariano, G. M., and Perinotti, P. (2011). ``Informational derivation of quantum theory.'' \emph{Physical Review A} \textbf{84}, 012311. \url{https://doi.org/10.1103/PhysRevA.84.012311}

\bibitem{pwchitambar2019quantum}
Chitambar, E., and Gour, G. (2019). ``Quantum resource theories.'' \emph{Reviews of Modern Physics} \textbf{91}, 025001. \url{https://doi.org/10.1103/RevModPhys.91.025001}

\bibitem{pwchristiansen2011discretizing}
Christiansen, S. H., and Halvorsen, T. G. (2011). ``Discretizing the Maxwell-Klein-Gordon equation by the lattice gauge theory formalism.'' \emph{IMA Journal of Numerical Analysis} \textbf{31}, 1--24. \url{https://doi.org/10.1093/imanum/drp019}

\bibitem{pwchristiansen2011gauge}
Christiansen, S. H., and Halvorsen, T. G. (2011). ``A gauge invariant discretization on simplicial grids of the Schr{\"{o}}dinger eigenvalue problem in an electromagnetic field.'' \emph{SIAM Journal on Numerical Analysis} \textbf{49}, 331--345. \url{https://doi.org/10.1137/090757502}

\bibitem{pwchristiansen2012simplicial}
Christiansen, S. H., and Halvorsen, T. G. (2012). ``A simplicial gauge theory.'' \emph{Journal of Mathematical Physics} \textbf{53}, 033501. \url{https://doi.org/10.1063/1.3692167}

\bibitem{pwchung1997spectral}
Chung, F. R. K. (1997). \emph{Spectral Graph Theory}. American Mathematical Society. \url{https://doi.org/10.1090/cbms/092}

\bibitem{pwchurch1936properties}
Church, A., and Rosser, J. B. (1936). ``Some properties of conversion.'' \emph{Transactions of the American Mathematical Society} \textbf{39}, 472--482. \url{https://doi.org/10.1090/S0002-9947-1936-1501858-0}

\bibitem{pwclauser1969proposed}
Clauser, J. F., Horne, M. A., Shimony, A., and Holt, R. A. (1969). ``Proposed experiment to test local hidden-variable theories.'' \emph{Physical Review Letters} \textbf{23}, 880--884. \url{https://doi.org/10.1103/PhysRevLett.23.880}

\bibitem{pwcombescure2012coherent}
Combescure, M., and Robert, D. (2012). \emph{Coherent States and Applications in Mathematical Physics}. Springer. \url{https://doi.org/10.1007/978-94-007-0196-0}

\bibitem{pwconnes1994neumann}
Connes, A., and Rovelli, C. (1994). ``Von Neumann algebra automorphisms and time-thermodynamics relation in generally covariant quantum theories.'' \emph{Classical and Quantum Gravity} \textbf{11}, 2899--2917. \url{https://doi.org/10.1088/0264-9381/11/12/007}

\bibitem{pwconway1985atlas}
Conway, J. H., Curtis, R. T., Norton, S. P., Parker, R. A., and Wilson, R. A. (1985). \emph{Atlas of Finite Groups: Maximal Subgroups and Ordinary Characters for Simple Groups}. Clarendon Press, Oxford. ISBN 0-19-853199-0.

\bibitem{pwcourant1928partiellen}
Courant, R., Friedrichs, K., and Lewy, H. (1928). ``{\"{U}}ber die partiellen Differenzengleichungen der mathematischen Physik.'' \emph{Mathematische Annalen} \textbf{100}, 32--74. \url{https://doi.org/10.1007/BF01448839}

\bibitem{pwcourant1967partial}
Courant, R., Friedrichs, K., and Lewy, H. (1967). ``On the partial difference equations of mathematical physics.'' \emph{IBM Journal of Research and Development} \textbf{11}, 215--234. \url{https://doi.org/10.1147/rd.112.0215}

\bibitem{pwcover2006elements}
Cover, T. M., and Thomas, J. A. (2006). \emph{Elements of Information Theory}. Wiley. \url{https://doi.org/10.1002/047174882X}

\bibitem{pwcoxeter1973regular}
Coxeter, H. S. M. (1973). \emph{Regular Polytopes}. Dover. ISBN 0-486-61480-8.

\bibitem{pwcurrie1966equivalent}
Currie, D. G., and Saletan, E. J. (1966). ``q-Equivalent Particle Hamiltonians. I. The Classical One-Dimensional Case.'' \emph{Journal of Mathematical Physics} \textbf{7}, 967--974. \url{https://doi.org/10.1063/1.1705010}

\bibitem{pwcurtis1962representation}
Curtis, C. W., and Reiner, I. (1962). \emph{Representation Theory of Finite Groups and Associative Algebras}. Interscience. \url{https://doi.org/10.1090/chel/356}

\bibitem{pwdaley2003introduction}
Daley, D. J., and Vere-Jones, D. (2003). \emph{An Introduction to the Theory of Point Processes, Volume I: Elementary Theory and Methods}. Springer, New York (Probability and its Applications). \url{https://doi.org/10.1007/b97277}

\bibitem{pwdavey2002introduction}
Davey, B. A., and Priestley, H. A. (2002). \emph{Introduction to Lattices and Order}. Cambridge University Press. \url{https://doi.org/10.1017/CBO9780511809088}

\bibitem{pwdavidson1996cstar}
Davidson, K. R. (1996). \emph{C*-Algebras by Example}. American Mathematical Society. \url{https://doi.org/10.1090/fim/006}

\bibitem{pwdavies1970operational}
Davies, E. B., and Lewis, J. T. (1970). ``An operational approach to quantum probability.'' \emph{Communications in Mathematical Physics} \textbf{17}, 239--260. \url{https://doi.org/10.1007/BF01647093}

\bibitem{pwdieks1982communication}
Dieks, D. (1982). ``Communication by EPR devices.'' \emph{Physics Letters A} \textbf{92}, 271--272. \url{https://doi.org/10.1016/0375-9601(82)90084-6}

\bibitem{pwdouglas1941solution}
Douglas, J. (1941). ``Solution of the inverse problem of the calculus of variations.'' \emph{Transactions of the American Mathematical Society} \textbf{50}, 71--128. \url{https://doi.org/10.1090/S0002-9947-1941-0004740-5}

\bibitem{pwdu1992analysis}
Du, Q., Gunzburger, M. D., and Peterson, J. S. (1992). ``Analysis and Approximation of the Ginzburg-Landau Model of Superconductivity.'' \emph{SIAM Review} \textbf{34}, 54--81. \url{https://doi.org/10.1137/1034003}

\bibitem{pweckmann1944harmonische}
Eckmann, B. (1944). ``Harmonische Funktionen und Randwertaufgaben in einem Komplex.'' \emph{Commentarii Mathematici Helvetici} \textbf{17}, 240--255. \url{https://doi.org/10.1007/BF02566245}

\bibitem{pwedelsbrunner2000edgewise}
Edelsbrunner, H., and Grayson, D. R. (2000). ``Edgewise subdivision of a simplex.'' \emph{Discrete \& Computational Geometry} \textbf{24}, 707--719. \url{https://doi.org/10.1007/s004540010063}

\bibitem{pweinstein1905elektrodynamik}
Einstein, A. (1905). ``Zur Elektrodynamik bewegter K{\"{o}}rper.'' \emph{Annalen der Physik} \textbf{322}, 891--921. \url{https://doi.org/10.1002/andp.19053221004}

\bibitem{pweinstein1905tragheit}
Einstein, A. (1905). ``Ist die Tr{\"{a}}gheit eines K{\"{o}}rpers von seinem Energieinhalt abh{\"{a}}ngig?.'' \emph{Annalen der Physik} \textbf{323}, 639--641. \url{https://doi.org/10.1002/andp.19053231314}

\bibitem{pweinstein1907uber}
Einstein, A. (1907). ``{\"{U}}ber die vom Relativit{\"{a}}tsprinzip geforderte Tr{\"{a}}gheit der Energie.'' \emph{Annalen der Physik} \textbf{328}, 371--384. \url{https://doi.org/10.1002/andp.19073280713}

\bibitem{pwevans2010partial}
Evans, L. C. (2010). \emph{Partial Differential Equations}. American Mathematical Society. \url{https://doi.org/10.1090/gsm/019}

\bibitem{pwfeller1968introduction}
Feller, W. (1968). \emph{An Introduction to Probability Theory and Its Applications, Volume 1}. Wiley. ISBN 0-471-25708-7.

\bibitem{pwfick1855diffusion}
Fick, A. (1855). ``Ueber Diffusion.'' \emph{Annalen der Physik} \textbf{170}, 59--86. \url{https://doi.org/10.1002/andp.18551700105}

\bibitem{pwfourier1822theorie}
Fourier, J. B. J. (1822). \emph{Th{\'{e}}orie analytique de la chaleur}. Firmin Didot, Paris; reissued by Cambridge University Press (Cambridge Library Collection), 2009. \url{https://doi.org/10.1017/CBO9780511693229}

\bibitem{pwfreudenthal1942simplizialzerlegungen}
Freudenthal, H. (1942). ``Simplizialzerlegungen von beschr{\"{a}}nkter Flachheit.'' \emph{Annals of Mathematics} \textbf{43}, 580. \url{https://doi.org/10.2307/1968813}

\bibitem{pwfulton1991representation}
Fulton, W., and Harris, J. (1991). \emph{Representation Theory: A First Course}. Springer, New York (Graduate Texts in Mathematics). \url{https://doi.org/10.1007/978-1-4612-0979-9}

\bibitem{pwgeman1984stochastic}
Geman, S., and Geman, D. (1984). ``Stochastic relaxation, Gibbs distributions, and the Bayesian restoration of images.'' \emph{IEEE Transactions on Pattern Analysis and Machine Intelligence} \textbf{PAMI-6}, 721--741. \url{https://doi.org/10.1109/TPAMI.1984.4767596}

\bibitem{pwgibbons1977cosmological}
Gibbons, G. W., and Hawking, S. W. (1977). ``Cosmological event horizons, thermodynamics, and particle creation.'' \emph{Physical Review D} \textbf{15}, 2738--2751. \url{https://doi.org/10.1103/PhysRevD.15.2738}

\bibitem{pwgleason1957measures}
Gleason, A. M. (1957). ``Measures on the closed subspaces of a Hilbert space.'' \emph{Journal of Mathematics and Mechanics} \textbf{6}, 885--893. \url{https://doi.org/10.1512/iumj.1957.6.56050}

\bibitem{pwgorini1976completely}
Gorini, V., Kossakowski, A., and Sudarshan, E. C. G. (1976). ``Completely positive dynamical semigroups of N-level systems.'' \emph{Journal of Mathematical Physics} \textbf{17}, 821--825. \url{https://doi.org/10.1063/1.522979}

\bibitem{pwgrady2010discrete}
Grady, L. J., and Polimeni, J. R. (2010). \emph{Discrete Calculus}. Springer. \url{https://doi.org/10.1007/978-1-84996-290-2}

\bibitem{pwgronwall1919note}
Gronwall, T. H. (1919). ``Note on the derivatives with respect to a parameter of the solutions of a system of differential equations.'' \emph{Annals of Mathematics} \textbf{20}, 292. \url{https://doi.org/10.2307/1967124}

\bibitem{pwhaag1964algebraic}
Haag, R., and Kastler, D. (1964). ``An algebraic approach to quantum field theory.'' \emph{Journal of Mathematical Physics} \textbf{5}, 848--861. \url{https://doi.org/10.1063/1.1704187}

\bibitem{pwhaag1996local}
Haag, R. (1996). \emph{Local Quantum Physics}. Springer. \url{https://doi.org/10.1007/978-3-642-61458-3}

\bibitem{pwhairer1993solving}
Hairer, E., N{\o}rsett, S. P., and Wanner, G. (1993). \emph{Solving Ordinary Differential Equations I: Nonstiff Problems}. Springer, Berlin (Springer Series in Computational Mathematics). \url{https://doi.org/10.1007/978-3-540-78862-1}

\bibitem{pwhairer2006geometric}
Hairer, E., Lubich, C., and Wanner, G. (2006). \emph{Geometric Numerical Integration: Structure-Preserving Algorithms for Ordinary Differential Equations}. Springer. \url{https://doi.org/10.1007/3-540-30666-8}

\bibitem{pwhartman2002ordinary}
Hartman, P. (2002). \emph{Ordinary Differential Equations}. Society for Industrial and Applied Mathematics (Classics in Applied Mathematics). \url{https://doi.org/10.1137/1.9780898719222}

\bibitem{pwheinosaari2011mathematical}
Heinosaari, T., and Ziman, M. (2011). \emph{The Mathematical Language of Quantum Theory: From Uncertainty to Entanglement}. Cambridge University Press. \url{https://doi.org/10.1017/CBO9781139031103}

\bibitem{pwhelmholtz1887physikalische}
von Helmholtz, H. (1887). ``Ueber die physikalische Bedeutung des Princips der kleinsten Wirkung.'' \emph{Journal f{\"{u}}r die reine und angewandte Mathematik} \textbf{100}, 137--166. \url{https://doi.org/10.1515/crll.1887.100.137}

\bibitem{pwhenneaux1982equations}
Henneaux, M. (1982). ``Equations of motion, commutation relations and ambiguities in the Lagrangian formalism.'' \emph{Annals of Physics} \textbf{140}, 45--64. \url{https://doi.org/10.1016/0003-4916(82)90334-7}

\bibitem{pwhenneaux1992quantization}
Henneaux, M., and Teitelboim, C. (1992). \emph{Quantization of Gauge Systems}. Princeton University Press. \url{https://doi.org/10.1515/9780691213866}

\bibitem{pwhodges1993model}
Hodges, W. (1993). \emph{Model Theory}. Cambridge University Press. \url{https://doi.org/10.1017/CBO9780511551574}

\bibitem{pwhydon2004variational}
Hydon, P. E., and Mansfield, E. L. (2004). ``A Variational Complex for Difference Equations.'' \emph{Foundations of Computational Mathematics} \textbf{4}, 187--217. \url{https://doi.org/10.1007/s10208-002-0071-9}

\bibitem{pwisaacs1976character}
Isaacs, I. M. (1976). \emph{Character Theory of Finite Groups}. Academic Press. \url{https://doi.org/10.1090/chel/359}

\bibitem{pwjackson1999classical}
Jackson, J. D. (1999). \emph{Classical Electrodynamics}. Wiley, New York. ISBN 0-471-30932-X.

\bibitem{pwjacobi1884vorlesungen}
Jacobi, C. G. J. (1884). \emph{Vorlesungen {\"{u}}ber Dynamik}. G. Reimer, Berlin. \url{https://doi.org/10.5962/bhl.title.18726}

\bibitem{pwjacobs1999categorical}
Jacobs, B. (1999). \emph{Categorical Logic and Type Theory}. Elsevier. \url{https://doi.org/10.1016/S0049-237X(98)X8028-6}

\bibitem{pwjames2001representations}
James, G., and Liebeck, M. (2001). \emph{Representations and Characters of Groups}. Cambridge University Press. \url{https://doi.org/10.1017/CBO9780511814532}

\bibitem{pwjaynes1957information}
Jaynes, E. T. (1957). ``Information Theory and Statistical Mechanics.'' \emph{Physical Review} \textbf{106}, 620--630. \url{https://doi.org/10.1103/PhysRev.106.620}

\bibitem{pwjohnson1994continuous}
Johnson, N. L., Kotz, S., and Balakrishnan, N. (1994). \emph{Continuous Univariate Distributions, Volume 1}. Wiley. ISBN 0-471-58495-9.

\bibitem{pwkac1947notion}
Kac, M. (1947). ``On the notion of recurrence in discrete stochastic processes.'' \emph{Bulletin of the American Mathematical Society} \textbf{53}, 1002--1010. \url{https://doi.org/10.1090/S0002-9904-1947-08927-8}

\bibitem{pwkato1995perturbation}
Kato, T. (1995). \emph{Perturbation Theory for Linear Operators}. Springer, Berlin (Classics in Mathematics). \url{https://doi.org/10.1007/978-3-642-66282-9}

\bibitem{pwkemeny1976finite}
Kemeny, J. G., and Snell, J. L. (1976). \emph{Finite Markov Chains}. Springer-Verlag, New York (Undergraduate Texts in Mathematics). ISBN 0-387-90192-2.

\bibitem{pwkilp2000monoids}
Kilp, M., Knauer, U., and Mikhalev, A. V. (2000). \emph{Monoids, Acts and Categories}. De Gruyter. \url{https://doi.org/10.1515/9783110812909}

\bibitem{pwkogut1979introduction}
Kogut, J. B. (1979). ``An introduction to lattice gauge theory and spin systems.'' \emph{Reviews of Modern Physics} \textbf{51}, 659--713. \url{https://doi.org/10.1103/RevModPhys.51.659}

\bibitem{pwkomorowski2012fluctuations}
Komorowski, T., Landim, C., and Olla, S. (2012). \emph{Fluctuations in Markov Processes: Time Symmetry and Martingale Approximation}. Springer, Berlin (Grundlehren der mathematischen Wissenschaften). \url{https://doi.org/10.1007/978-3-642-29880-6}

\bibitem{pwkorner1998zeroerror}
K{\"{o}}rner, J., and Orlitsky, A. (1998). ``Zero-error information theory.'' \emph{IEEE Transactions on Information Theory} \textbf{44}, 2207--2229. \url{https://doi.org/10.1109/18.720537}

\bibitem{pwkraus1983states}
Kraus, K. (1983). \emph{States, Effects, and Operations: Fundamental Notions of Quantum Theory}. Springer. \url{https://doi.org/10.1007/3-540-12732-1}

\bibitem{pwkronheimer1967structure}
Kronheimer, E. H., and Penrose, R. (1967). ``On the structure of causal spaces.'' \emph{Proceedings of the Cambridge Philosophical Society} \textbf{63}, 481--501. \url{https://doi.org/10.1017/S030500410004144X}

\bibitem{pwkuhn1960some}
Kuhn, H. W. (1960). ``Some combinatorial lemmas in topology.'' \emph{IBM Journal of Research and Development} \textbf{4}, 518--524. \url{https://doi.org/10.1147/rd.45.0518}

\bibitem{pwlam1998theorem}
Lam, T. Y. (1998). ``A theorem of Burnside on matrix rings.'' \emph{American Mathematical Monthly} \textbf{105}, 651--653. \url{https://doi.org/10.1080/00029890.1998.12004939}

\bibitem{pwlamport1978time}
Lamport, L. (1978). ``Time, clocks, and the ordering of events in a distributed system.'' \emph{Communications of the ACM} \textbf{21}, 558--565. \url{https://doi.org/10.1145/359545.359563}

\bibitem{pwlandau1975classical}
Landau, L. D., and Lifshitz, E. M. (1975). \emph{The Classical Theory of Fields}. Pergamon Press. ISBN 978-0-08-025072-4.

\bibitem{pwlaurie2013certificate}
Laurie, B., Langley, A., and Kasper, E. (2013). ``Certificate Transparency.'' \url{https://doi.org/10.17487/RFC6962}

\bibitem{pwleblanc1987debugging}
LeBlanc, T. J., and Mellor-Crummey, J. M. (1987). ``Debugging parallel programs with Instant Replay.'' \emph{IEEE Transactions on Computers} \textbf{C-36}, 471--482. \url{https://doi.org/10.1109/TC.1987.1676929}

\bibitem{pwlehmann1998theory}
Lehmann, E. L., and Casella, G. (1998). \emph{Theory of Point Estimation}. Springer, New York (Springer Texts in Statistics). \url{https://doi.org/10.1007/b98854}

\bibitem{pwleimkuhler2004simulating}
Leimkuhler, B., and Reich, S. (2004). \emph{Simulating Hamiltonian Dynamics}. Cambridge University Press. \url{https://doi.org/10.1017/CBO9780511614118}

\bibitem{pwlevin2017markov}
Levin, D. A., and Peres, Y. (2017). \emph{Markov Chains and Mixing Times}. American Mathematical Society. \url{https://doi.org/10.1090/mbk/107}

\bibitem{pwlim2020hodge}
Lim, L. H. (2020). ``Hodge Laplacians on graphs.'' \emph{SIAM Review} \textbf{62}, 685--715. \url{https://doi.org/10.1137/18M1223101}

\bibitem{pwlindblad1975completely}
Lindblad, G. (1975). ``Completely positive maps and entropy inequalities.'' \emph{Communications in Mathematical Physics} \textbf{40}, 147--151. \url{https://doi.org/10.1007/BF01609396}

\bibitem{pwlindblad1976generators}
Lindblad, G. (1976). ``On the generators of quantum dynamical semigroups.'' \emph{Communications in Mathematical Physics} \textbf{48}, 119--130. \url{https://doi.org/10.1007/BF01608499}

\bibitem{pwludascher2006scientific}
Lud{\"{a}}scher, B., Altintas, I., Berkley, C., Higgins, D., Jaeger, E., Jones, M., et al. (2006). ``Scientific workflow management and the Kepler system.'' \emph{Concurrency and Computation: Practice and Experience} \textbf{18}, 1039--1065. \url{https://doi.org/10.1002/cpe.994}

\bibitem{pwluders1950zustandsanderung}
L{\"{u}}ders, G. (1950). ``{\"{U}}ber die Zustands{\"{a}}nderung durch den Me{\ss}proze{\ss}.'' \emph{Annalen der Physik} \textbf{443}, 322--328. \url{https://doi.org/10.1002/andp.19504430510}

\bibitem{pwmaclane1992sheaves}
Mac Lane, S., and Moerdijk, I. (1992). \emph{Sheaves in Geometry and Logic: A First Introduction to Topos Theory}. Springer. \url{https://doi.org/10.1007/978-1-4612-0927-0}

\bibitem{pwmaclane1998categories}
Mac Lane, S. (1998). \emph{Categories for the Working Mathematician}. Springer. \url{https://doi.org/10.1007/978-1-4757-4721-8}

\bibitem{pwmandula1983representations}
Mandula, J. E., Zweig, G., and Govaerts, J. (1983). ``Representations of the rotation reflection symmetry group of the four-dimensional cubic lattice.'' \emph{Nuclear Physics B} \textbf{228}, 91--108. \url{https://doi.org/10.1016/0550-3213(83)90399-1}

\bibitem{pwmarsden1974reduction}
Marsden, J., and Weinstein, A. (1974). ``Reduction of symplectic manifolds with symmetry.'' \emph{Reports on Mathematical Physics} \textbf{5}, 121--130. \url{https://doi.org/10.1016/0034-4877(74)90021-4}

\bibitem{pwmarsden1999introduction}
Marsden, J. E., and Ratiu, T. S. (1999). \emph{Introduction to Mechanics and Symmetry: A Basic Exposition of Classical Mechanical Systems}. Springer, New York (Texts in Applied Mathematics). \url{https://doi.org/10.1007/978-0-387-21792-5}

\bibitem{pwmasanes2017general}
Masanes, L., and Oppenheim, J. (2017). ``A general derivation and quantification of the third law of thermodynamics.'' \emph{Nature Communications} \textbf{8}, 14538. \url{https://doi.org/10.1038/ncomms14538}

\bibitem{pwmcduff2017introduction}
McDuff, D., and Salamon, D. (2017). \emph{Introduction to Symplectic Topology}. Oxford University Press. \url{https://doi.org/10.1093/oso/9780198794899.001.0001}

\bibitem{pwmontvay1994quantum}
Montvay, I., and M{\"{u}}nster, G. (1994). \emph{Quantum Fields on a Lattice}. Cambridge University Press. \url{https://doi.org/10.1017/CBO9780511470783}

\bibitem{pwmoore2009introduction}
Moore, R. E., Kearfott, R. B., and Cloud, M. J. (2009). \emph{Introduction to Interval Analysis}. SIAM. \url{https://doi.org/10.1137/1.9780898717716}

\bibitem{pwnaber2012geometry}
Naber, G. L. (2012). \emph{The Geometry of Minkowski Spacetime: An Introduction to the Mathematics of the Special Theory of Relativity}. Springer. \url{https://doi.org/10.1007/978-1-4419-7838-7}

\bibitem{pwnedelec1980mixed}
N{\'{e}}d{\'{e}}lec, J. C. (1980). ``Mixed finite elements in $\mathbf{R}^3$.'' \emph{Numerische Mathematik} \textbf{35}, 315--341. \url{https://doi.org/10.1007/BF01396415}

\bibitem{pwnedialkov1999validated}
Nedialkov, N. S., Jackson, K. R., and Corliss, G. F. (1999). ``Validated solutions of initial value problems for ordinary differential equations.'' \emph{Applied Mathematics and Computation} \textbf{105}, 21--68. \url{https://doi.org/10.1016/S0096-3003(98)10083-8}

\bibitem{pwnewman1942theories}
Newman, M. H. A. (1942). ``On theories with a combinatorial definition of "equivalence".'' \emph{Annals of Mathematics} \textbf{43}, 223. \url{https://doi.org/10.2307/1968867}

\bibitem{pwnielsen2010quantum}
Nielsen, M. A., and Chuang, I. L. (2010). \emph{Quantum Computation and Quantum Information: 10th Anniversary Edition}. Cambridge University Press. \url{https://doi.org/10.1017/CBO9780511976667}

\bibitem{pwohya1993quantum}
Ohya, M., and Petz, D. (1993). \emph{Quantum Entropy and Its Use}. Springer. \url{https://doi.org/10.1007/978-3-642-57997-4}

\bibitem{pwolver1993applications}
Olver, P. J. (1993). \emph{Applications of Lie Groups to Differential Equations}. Springer, New York (Graduate Texts in Mathematics). \url{https://doi.org/10.1007/978-1-4612-4350-2}

\bibitem{pwolver2010nist}
Olver, F. W. J., Lozier, D. W., Boisvert, R. F., and Clark, C. W. (2010). \emph{NIST Handbook of Mathematical Functions}. Cambridge University Press. ISBN 978-0-521-19225-5.

\bibitem{pwoneill1983semi}
O'Neill, B. (1983). \emph{Semi-Riemannian Geometry with Applications to Relativity}. Academic Press, New York. ISBN 0-12-526740-1.

\bibitem{pwpaulsen2002completely}
Paulsen, V. (2002). \emph{Completely Bounded Maps and Operator Algebras}. Cambridge University Press. \url{https://doi.org/10.1017/CBO9780511546631}

\bibitem{pwpazy1983semigroups}
Pazy, A. (1983). \emph{Semigroups of Linear Operators and Applications to Partial Differential Equations}. Springer. \url{https://doi.org/10.1007/978-1-4612-5561-1}

\bibitem{pwpenrose1984spinors}
Penrose, R., and Rindler, W. (1984). \emph{Spinors and Space-Time, Volume 1}. Cambridge University Press. \url{https://doi.org/10.1017/CBO9780511564048}

\bibitem{pwpezoa2016foundations}
Pezoa, F., Reutter, J. L., Suarez, F., Ugarte, M., and Vrgo{\v{c}}, D. (2016). ``Foundations of JSON Schema.'' In \emph{Proceedings of the 25th International Conference on World Wide Web (WWW '16)}, 263--273. International World Wide Web Conferences Steering Committee. \url{https://doi.org/10.1145/2872427.2883029}

\bibitem{pwpierce2002types}
Pierce, B. C. (2002). \emph{Types and Programming Languages}. MIT Press. ISBN 0-262-16209-1.

\bibitem{pwplotkin2004structural}
Plotkin, G. D. (2004). ``A structural approach to operational semantics.'' \emph{The Journal of Logic and Algebraic Programming} \textbf{60-61}, 17--139. \url{https://doi.org/10.1016/j.jlap.2004.05.001}

\bibitem{pwpolchinski1998string}
Polchinski, J. (1998). \emph{String Theory. Volume 1}. Cambridge University Press. \url{https://doi.org/10.1017/CBO9780511816079}

\bibitem{pwprugovecki1977information}
Prugove{\v{c}}ki, E. (1977). ``Information-theoretical aspects of quantum measurement.'' \emph{International Journal of Theoretical Physics} \textbf{16}, 321--331. \url{https://doi.org/10.1007/BF01807146}

\bibitem{pwratcliffe2019foundations}
Ratcliffe, J. G. (2019). \emph{Foundations of Hyperbolic Manifolds}. Springer International Publishing. \url{https://doi.org/10.1007/978-3-030-31597-9}

\bibitem{pwreed1980methods}
Reed, M., and Simon, B. (1980). \emph{Methods of Modern Mathematical Physics I: Functional Analysis}. Academic Press. ISBN 0-12-585050-6.

\bibitem{pwriehle2004frequency}
Riehle, F. (2004). \emph{Frequency Standards: Basics and Applications}. Wiley-VCH, Weinheim. \url{https://doi.org/10.1002/3527605991}

\bibitem{pwrindler2006relativity}
Rindler, W. (2006). \emph{Relativity: Special, General, and Cosmological}. Oxford University Press. \url{https://doi.org/10.1093/oso/9780198567318.001.0001}

\bibitem{pwringer2019qed}
Ringer, T., Palmskog, K., Sergey, I., Gligoric, M., and Tatlock, Z. (2019). ``QED at Large: A Survey of Engineering of Formally Verified Software.'' \emph{Foundations and Trends in Programming Languages} \textbf{5}, 102--281. \url{https://doi.org/10.1561/2500000045}

\bibitem{pwrockafellar1970convex}
Rockafellar, R. T. (1970). \emph{Convex Analysis}. Princeton University Press. \url{https://doi.org/10.1515/9781400873173}

\bibitem{pwrogers1981markov}
Rogers, L. C. G., and Pitman, J. W. (1981). ``Markov Functions.'' \emph{The Annals of Probability} \textbf{9}, 573--582. \url{https://doi.org/10.1214/aop/1176994363}

\bibitem{pwrudin1991functional}
Rudin, W. (1991). \emph{Functional Analysis}. McGraw-Hill. ISBN 0-07-054236-8.

\bibitem{pwrundgren2020json}
Rundgren, A., Jordan, B., and Erdtman, S. (2020). ``JSON Canonicalization Scheme (JCS).'' \url{https://doi.org/10.17487/RFC8785}

\bibitem{pwsandve2013ten}
Sandve, G. K., Nekrutenko, A., Taylor, J., and Hovig, E. (2013). ``Ten simple rules for reproducible computational research.'' \emph{PLoS Computational Biology} \textbf{9}, e1003285. \url{https://doi.org/10.1371/journal.pcbi.1003285}

\bibitem{pwschur1905neue}
Schur, I. (1905). ``Neue Begr{\"{u}}ndung der Theorie der Gruppencharaktere.'' \emph{Sitzungsberichte der K{\"{o}}niglich Preu{\ss}ischen Akademie der Wissenschaften zu Berlin}, 406--432. \url{https://zbmath.org/?q=an:36.0194.01}

\bibitem{pwsegal1963global}
Segal, I. E. (1963). ``The global Cauchy problem for a relativistic scalar field with power interaction.'' \emph{Bulletin de la Soci{\'{e}}t{\'{e}} Math{\'{e}}matique de France} \textbf{91}, 129--135. \url{https://doi.org/10.24033/bsmf.1593}

\bibitem{pwseifert2012stochastic}
Seifert, U. (2012). ``Stochastic thermodynamics, fluctuation theorems and molecular machines.'' \emph{Reports on Progress in Physics} \textbf{75}, 126001. \url{https://doi.org/10.1088/0034-4885/75/12/126001}

\bibitem{pwserre1977linear}
Serre, J. P. (1977). \emph{Linear Representations of Finite Groups}. Springer. \url{https://doi.org/10.1007/978-1-4684-9458-7}

\bibitem{pwshannon1956zero}
Shannon, C. (1956). ``The zero error capacity of a noisy channel.'' \emph{IRE Transactions on Information Theory} \textbf{2}, 8--19. \url{https://doi.org/10.1109/TIT.1956.1056798}

\bibitem{pwshatah2000geometric}
Shatah, J., and Struwe, M. (2000). \emph{Geometric Wave Equations}. American Mathematical Society. \url{https://doi.org/10.1090/cln/002}

\bibitem{pwsontag1998mathematical}
Sontag, E. D. (1998). \emph{Mathematical Control Theory}. Springer. \url{https://doi.org/10.1007/978-1-4612-0577-7}

\bibitem{pwspivak2014category}
Spivak, D. I. (2014). \emph{Category Theory for the Sciences}. The MIT Press. \url{https://doi.org/10.7551/mitpress/9995.001.0001}

\bibitem{pwstock2019revision}
Stock, M., Davis, R., de Mirand{\'{e}}s, E., and Milton, M. J. T. (2019). ``The revision of the SI---the result of three decades of progress in metrology.'' \emph{Metrologia} \textbf{56}, 022001. \url{https://doi.org/10.1088/1681-7575/ab0013}

\bibitem{pwstrauss1989nonlinear}
Strauss, W. A. (1989). \emph{Nonlinear Wave Equations}. American Mathematical Society (CBMS Regional Conference Series in Mathematics 73). \url{https://doi.org/10.1090/cbms/073}

\bibitem{pwsummers1990independence}
Summers, S. J. (1990). ``On the independence of local algebras in quantum field theory.'' \emph{Reviews in Mathematical Physics} \textbf{2}, 201--247. \url{https://doi.org/10.1142/S0129055X90000090}

\bibitem{pwtakesaki1972conditional}
Takesaki, M. (1972). ``Conditional expectations in von Neumann algebras.'' \emph{Journal of Functional Analysis} \textbf{9}, 306--321. \url{https://doi.org/10.1016/0022-1236(72)90004-3}

\bibitem{pwtarantola2005inverse}
Tarantola, A. (2005). \emph{Inverse Problem Theory and Methods for Model Parameter Estimation}. SIAM. \url{https://doi.org/10.1137/1.9780898717921}

\bibitem{pwtarski1955lattice}
Tarski, A. (1955). ``A lattice-theoretical fixpoint theorem and its applications.'' \emph{Pacific Journal of Mathematics} \textbf{5}, 285--309. \url{https://doi.org/10.2140/pjm.1955.5.285}

\bibitem{pwtomiyama1957projection}
Tomiyama, J. (1957). ``On the projection of norm one in W*-algebras.'' \emph{Proceedings of the Japan Academy} \textbf{33}, 608--612. \url{https://doi.org/10.3792/pja/1195524885}

\bibitem{pwuhlmann1977relative}
Uhlmann, A. (1977). ``Relative entropy and the Wigner-Yanase-Dyson-Lieb concavity in an interpolation theory.'' \emph{Communications in Mathematical Physics} \textbf{54}, 21--32. \url{https://doi.org/10.1007/BF01609834}

\bibitem{pwvonneumann1932mathematische}
von Neumann, J. (1932). \emph{Mathematische Grundlagen der Quantenmechanik}. Springer. \url{https://doi.org/10.1007/978-3-642-61409-5}

\bibitem{pwwatrous2018theory}
Watrous, J. (2018). \emph{The Theory of Quantum Information}. Cambridge University Press. \url{https://doi.org/10.1017/9781316848142}

\bibitem{pwwerner1989quantum}
Werner, R. F. (1989). ``Quantum states with Einstein-Podolsky-Rosen correlations admitting a hidden-variable model.'' \emph{Physical Review A} \textbf{40}, 4277--4281. \url{https://doi.org/10.1103/PhysRevA.40.4277}

\bibitem{pwwheeler1973priori}
Wheeler, M. F. (1973). ``A priori L2 error estimates for Galerkin approximations to parabolic partial differential equations.'' \emph{SIAM Journal on Numerical Analysis} \textbf{10}, 723--759. \url{https://doi.org/10.1137/0710062}

\bibitem{pwwhitney1957geometric}
Whitney, H. (1957). \emph{Geometric Integration Theory}. Princeton University Press. \url{https://doi.org/10.1515/9781400877577}

\bibitem{pwwilson1974confinement}
Wilson, K. G. (1974). ``Confinement of quarks.'' \emph{Physical Review D} \textbf{10}, 2445--2459. \url{https://doi.org/10.1103/PhysRevD.10.2445}

\bibitem{pwwilson1980molecular}
Wilson, J. E. B., Decius, J. C., and Cross, P. C. (1980). \emph{Molecular Vibrations: The Theory of Infrared and Raman Vibrational Spectra}. Dover Publications. ISBN 0-486-63941-X.

\bibitem{pwwootters1982single}
Wootters, W. K., and Zurek, W. H. (1982). ``A single quantum cannot be cloned.'' \emph{Nature} \textbf{299}, 802--803. \url{https://doi.org/10.1038/299802a0}

\bibitem{pwwreszinski2009precise}
Wreszinski, W. F., and Abdalla, E. (2009). ``A Precise Formulation of the Third Law of Thermodynamics.'' \emph{Journal of Statistical Physics} \textbf{134}, 781--792. \url{https://doi.org/10.1007/s10955-009-9693-5}

\bibitem{pwyee1966numerical}
Yee, K. S. (1966). ``Numerical solution of initial boundary value problems involving Maxwell's equations in isotropic media.'' \emph{IEEE Transactions on Antennas and Propagation} \textbf{14}, 302--307. \url{https://doi.org/10.1109/TAP.1966.1138693}

\bibitem{pwzeeman1964causality}
Zeeman, E. C. (1964). ``Causality implies the Lorentz group.'' \emph{Journal of Mathematical Physics} \textbf{5}, 490--493. \url{https://doi.org/10.1063/1.1704140}

